Given two vector spaces $V, W$ over the same field $F$, a mapping $T$ of $V$ into $W$ is called a homomorphism (or a linear transformation) if $(x + y)T = xT + yT$ for every $x$ and $y$ in $V$, and if $(ax)T = a(xT)$ for every $x$ in $V$ and every $a$ in $F$. Then $T$ is, in particular, a homomorphism of the additive group of $V$ into that of $W$.

The kernel of $T$ is a vector subspace of $V$, and the image $VT$ of $V$ is a vector subspace of $W$.

A homomorphism of $V$ into $W$ which is univalent (that is, whose kernel is $(0)$) is called an isomorphism of $V$ into $W$.

A homomorphism of $V$ into itself is called an endomorphism; an endomorphism of $V$ which is univalent and onto is called an automorphism of $V$.
